
\documentclass{article}
\usepackage{colortbl}
\usepackage{makecell}
\usepackage{multirow}
\usepackage{supertabular}

\begin{document}

\newcounter{utterance}

\centering \large Interaction Transcript for game `cladder', experiment `full\_v1.5\_default', episode 420 with Qwen3.6{-}27B.
\vspace{24pt}

{ \footnotesize  \setcounter{utterance}{1}
\setlength{\tabcolsep}{0pt}
\begin{supertabular}{c@{$\;$}|p{.15\linewidth}@{}p{.15\linewidth}p{.15\linewidth}p{.15\linewidth}p{.15\linewidth}p{.15\linewidth}}
    \# & \multicolumn{2}{c}{Player} && \multicolumn{2}{c}{Game Master} \\
    \hline

    \theutterance \stepcounter{utterance}  
    & & & \multicolumn{4}{p{0.6\linewidth}}{
        \cellcolor[rgb]{0.9,0.9,0.9}{
            \makecell[{{p{\linewidth}}}]{
                \texttt{\tiny{[P1$\langle$GM]}}
                \texttt{You are an expert in causal inference. The following question is not a typical commonsense query, but rather a meticulously designed question created by a professor specializing in causal inference, intended to assess the students' mastery of the course content.} \\
\texttt{Imagine a self{-}contained, hypothetical world with only the following conditions, and without any unmentioned factors or causal relationships: Vaccination status has a direct effect on vaccination reaction and getting smallpox. Getting smallpox has a direct effect on smallpox survival. Vaccination reaction has a direct effect on smallpox survival. Method 1: We look directly at how vaccination status correlates with smallpox survival in general. Method 2: We look at this correlation case by case according to vaccination reaction. To understand how vaccination status affects smallpox survival, is it more correct to use the Method 1 than Method 2?} \\
\texttt{Start your answer with 'Yes' or 'No', followed by additional reasoning or evidence to support your explanation.} \\
            }
        }
    }
     \\ \\

    \theutterance \stepcounter{utterance}  
    & \multicolumn{4}{p{0.6\linewidth}}{
        \cellcolor[rgb]{0.561,0.737,0.561}{
            \makecell[{{p{\linewidth}}}]{
                \texttt{\tiny{[P1$\rangle$GM]}}
                \texttt{Yes} \\
\\ 
\texttt{**Reasoning:**} \\
\\ 
\texttt{To determine the correct method, we must analyze the causal structure described in the hypothetical world using a Directed Acyclic Graph (DAG).} \\
\\ 
\texttt{1.  **Identify the Variables and Causal Links:**} \\
\texttt{    *   \$V\$: Vaccination Status} \\
\texttt{    *   \$R\$: Vaccination Reaction} \\
\texttt{    *   \$S\$: Getting Smallpox} \\
\texttt{    *   \$Y\$: Smallpox Survival} \\
\\ 
\texttt{    The problem states the following direct effects (edges in the DAG):} \\
\texttt{    *   \$V \textbackslash{}rightarrow R\$} \\
\texttt{    *   \$V \textbackslash{}rightarrow S\$} \\
\texttt{    *   \$S \textbackslash{}rightarrow Y\$} \\
\texttt{    *   \$R \textbackslash{}rightarrow Y\$} \\
\\ 
\texttt{    There are no other unmentioned factors or relationships.} \\
\\ 
\texttt{2.  **Identify the Causal Paths from \$V\$ to \$Y\$:**} \\
\texttt{    We want to estimate the total causal effect of Vaccination Status (\$V\$) on Smallpox Survival (\$Y\$). Let's trace the paths from \$V\$ to \$Y\$:} \\
\texttt{    *   **Path 1:** \$V \textbackslash{}rightarrow S \textbackslash{}rightarrow Y\$ (This is a causal path representing the effect of vaccination on survival via preventing/causing smallpox).} \\
\texttt{    *   **Path 2:** \$V \textbackslash{}rightarrow R \textbackslash{}rightarrow Y\$ (This is a causal path representing the effect of vaccination on survival via the vaccination reaction).} \\
\\ 
\texttt{    Both paths are **causal} \\
            }
        }
    }
    & & \\ \\

    \theutterance \stepcounter{utterance}  
    & & & \multicolumn{2}{p{0.3\linewidth}}{
        \cellcolor[rgb]{0.9,0.9,0.9}{
            \makecell[{{p{\linewidth}}}]{
                \texttt{\tiny{[GM$|$GM]}}
                \texttt{Yes} \\
            }
        }
    }
    & & \\ \\

    \theutterance \stepcounter{utterance}  
    & & & \multicolumn{2}{p{0.3\linewidth}}{
        \cellcolor[rgb]{0.9,0.9,0.9}{
            \makecell[{{p{\linewidth}}}]{
                \texttt{\tiny{[GM$|$GM]}}
                \texttt{yes} \\
            }
        }
    }
    & & \\ \\

    \theutterance \stepcounter{utterance}  
    & & & \multicolumn{2}{p{0.3\linewidth}}{
        \cellcolor[rgb]{0.9,0.9,0.9}{
            \makecell[{{p{\linewidth}}}]{
                \texttt{\tiny{[GM$|$GM]}}
                \texttt{game\_result = WIN} \\
            }
        }
    }
    & & \\ \\

\end{supertabular}
}

\end{document}
